\documentclass{article}
\usepackage{mgvalues}
\title{Reporting of hepatic injury with example-drug: a disproportionality analysis of a synthetic spontaneous reporting database}

\begin{document}
\maketitle

\begin{abstract}
\textbf{Background.} Drug-induced hepatic injury is a common reason for regulatory action.
\textbf{Methods.} Disproportionality analysis of a synthetic spontaneous reporting database
covering \mg{results.cohort.period_start} to \mg{results.cohort.period_end}, comparing
reports of hepatic injury for example-drug against all other drugs.
\textbf{Results.} Of \mg{results.cohort.n_reports} reports, \mg{results.case.n_cases} described
hepatic injury with example-drug. The reporting odds ratio was \mg{results.ror.point}
(95\% CI \mg{results.ror.ci_low} to \mg{results.ror.ci_high}).
\textbf{Conclusions.} Reporting of hepatic injury is disproportionate for example-drug in these
synthetic data. Disproportionality is not incidence and cannot establish risk.
\end{abstract}

\section{Introduction}

Drug-induced hepatic injury remains among the commonest reasons for regulatory action on a
marketed medicine. In a national cohort, the crude incidence was
\mg{lit.background.reported_incidence} per 100\,000 person-years
\cite{fictionalHepaticCohort2021}, and a disproportionality signal has been reported for the
drug class as a whole (ROR \mg{lit.background.class_ror}) \cite{fictionalClassSignal2019}. The
national agency estimated that \mg{lit.agency.withdrawn_estimate} patients had been exposed
by the end of its last published review.

Whether the signal extends to example-drug specifically has not been examined.

\section{Methods}

We analysed a synthetic spontaneous reporting database, generated on
\mg{results.dataset.reports.extracted} to cover \mg{results.dataset.reports.covers_from} to
\mg{results.dataset.reports.covers_to}. Its reports run from
\mg{results.cohort.period_start} to \mg{results.cohort.period_end}, a period of
\mg{results.cohort.n_years} years. Reports were included without restriction on age or sex.
Each record carries a unique report identifier; checked against the data dictionary, \mg{results.quality.rule.one_row_per_report} of the \mg{results.quality.rows}
reports repeated another's and \mg{results.quality.age_group.missing} gave no age group. The generator emits no duplicate cases, so no de-duplication
step was applied; a real database would require one, and its absence here is a property of the synthetic data rather than a design choice.

Hepatic injury is the single event term used by the data generator. Real work would need a
case definition --- a MedDRA Standardised Query or an explicit preferred-term list --- and the
choice of definition would change the counts. The terms actually used are published in
Table~S1, emitted from the lists the analysis filtered on rather than retyped, which is what
RECORD 6.1 asks for and what makes a case definition checkable.

The reporting odds ratio was computed from a 2 x 2 table contrasting reports of hepatic
injury with all other reported events, for example-drug against all other drugs in the
database. Confidence intervals were derived from the standard error of the log odds ratio
and are 95\% confidence intervals, two-sided at an alpha of \mg{results.param.alpha}; the
main estimate is also given with its 90\% confidence interval, for comparison with
signal-detection practice. \mg{method.ror} \mg{method.ci}

A signal was defined by the classical criterion: at least
\mg{results.param.signal.min_cases} cases together with a lower bound of the 95\% confidence
interval of the reporting odds ratio above 1. No p-value threshold was used as a decision
rule, which is not how a disproportionality analysis reaches its conclusion.
\mg{method.signal}

Adjusted for age group and sex, the two patient characteristics the database records for
every report, the reporting odds ratio was also estimated by \mg{results.model.adjusted.kind}
of hepatic injury on exposure to example-drug, with both characteristics entered as
categories. The reporting year, which every report also carries, was not adjusted for:
calendar time is a common adjustment in disproportionality analyses, but the data generator
draws the year independently of the drug, so here it cannot confound the estimate. Its
interval is the model's Wald interval for the exposure coefficient, at the same alpha. Table
S2 describes the model as it was fitted. \mg{method.adjusted}

The analysis was run in Python \mg{results.software.python} with statsmodels
\mg{results.software.statsmodels}.

Reporting follows the checklist declared in \texttt{paper.yaml}.

% That checklist is a demonstration profile, not a published guideline, which is why no
% guideline is named in the prose above. As in the example's main.md, the comment is the
% twin's too, and its numbers, RECORD-PE and 2019, are read by no gate.

\section{Results}

The database contained \mg{results.cohort.n_reports} reports, of which
\mg{results.cohort.n_drug_reports} named example-drug. Hepatic injury was reported in
\mg{results.case.n_cases} of these; \mg{results.case.n_serious} were flagged as serious
(\mg{results.case.pct_serious}\%). Baseline characteristics are shown in Table~1.

\mg{table.baseline}

Reporting of hepatic injury was disproportionate for example-drug, with a reporting odds
ratio of \mg{results.ror.point} (95\% CI \mg{results.ror.ci_low} to
\mg{results.ror.ci_high}; 90\% CI \mg{results.ror.ci90_low} to \mg{results.ror.ci90_high}). The
signal criterion defined in the Methods was \mg{results.signal.verdict}. The estimate is shown
in Figure~1, and the counts it was computed from in Table~2, so a reader
can reconstruct it. Adjusted for age group and sex, the reporting odds ratio was
\mg{results.ror_adjusted.point} (95\% CI \mg{results.ror_adjusted.ci_low} to
\mg{results.ror_adjusted.ci_high}), given to three decimal places because to two it is the
crude estimate: the data generator draws both characteristics independently of the drug.

\mg{table.two_by_two}

\mg{figure.forest}

\section{Discussion}

The reporting odds ratio observed here exceeds the class-level estimate of
\mg{lit.background.class_ror} reported previously \cite{fictionalClassSignal2019}, although the
two are not directly comparable: that analysis pooled the class and drew on a different
reporting period.

Disproportionality is not incidence, and nothing in these data supports a rate. The cohort
incidence of \mg{lit.background.reported_incidence} per 100\,000 person-years
\cite{fictionalHepaticCohort2021} is the only population-level figure available, and it
predates the marketing of example-drug.

Several limitations follow from the design. Spontaneous reports are subject to notoriety
bias, and the denominator is unknown. Grade 3 and higher events cannot be distinguished
from milder ones in these data. The exposure figure of
\mg{lit.agency.withdrawn_estimate} patients rests on a printed agency report that is no
longer publicly available, and is recorded as an author attestation rather than a stored
source.

\section{Data availability}

The synthetic dataset, the code that generates it and the analysis that reads it are all in
this repository. No real patient data were used, and no ethical approval was required.

\section{Funding}

This work received no funding.

\section{Competing interests}

None declared.

\bibliographystyle{plain}
\bibliography{references}
\end{document}
